\documentclass[12pt,a4paper]{article}

\usepackage{cmap}
\usepackage[T2A]{fontenc}
\usepackage[utf8]{inputenc}
\usepackage[russian]{babel}
\usepackage{amsmath,amssymb}
\usepackage[a4paper,margin=2cm]{geometry}

\newcommand{\R}{\mathbb{R}}
\newcommand{\eps}{\varepsilon}
\newcommand{\Int}{\operatorname{int}}
\newcommand{\Fr}{\partial}
\newcommand{\lowmu}{\underline{\mu}}
\newcommand{\upmu}{\overline{\mu}}

\setlength{\parindent}{0pt}
\setlength{\parskip}{0.45em}
\sloppy

\begin{document}

\begin{center}
{\Large\bfseries Билет 5}
\end{center}

\noindent\fbox{%
  \begin{minipage}{\dimexpr\linewidth-2\fboxsep-2\fboxrule\relax}
  \normalfont
Определение измеримости по Жордану множества в \(n\)-мерном евклидовом
пространстве. Критерий измеримости. Измеримость объединения, пересечения и
разности измеримых множеств. Конечная аддитивность меры Жордана.
  \end{minipage}%
}

\section*{Краткий план}

Сначала вводятся клетка, клеточное множество и его мера. Затем через
приближение множества изнутри и снаружи клеточными множествами определяются
нижняя и верхняя меры Жордана; множество измеримо, когда эти меры равны и
конечны. Основной критерий: множество в \(\R^n\) измеримо по Жордану тогда и
только тогда, когда оно ограничено и его граница имеет меру нуль. Из критерия
следует измеримость объединения, пересечения и разности измеримых множеств,
потому что границы этих множеств лежат в объединении исходных границ. В конце
доказывается конечная аддитивность: для измеримых множеств без общих
внутренних точек мера объединения равна сумме мер.

\section*{Определения}

\textbf{Клетка.}
Клеткой в \(\R^n\) называется замкнутый прямоугольный параллелепипед
\[
  \Pi=
  \{x=(x_1,\dots,x_n): a_k\le x_k\le b_k,\ k=1,\dots,n\},
\]
где \(a_k\le b_k\). Его мерой называется
\[
  \mu(\Pi)=\prod_{k=1}^n(b_k-a_k).
\]

\textbf{Клеточное множество.}
Множество \(A\subset\R^n\) называется клеточным, если оно представимо в виде
конечного объединения клеток с попарно непересекающимися внутренностями:
\[
  A=\bigcup_{i=1}^N \Pi_i,\qquad
  \Int\Pi_i\cap\Int\Pi_j=\varnothing\quad(i\ne j).
\]
Его мера равна
\[
  \mu(A)=\sum_{i=1}^N\mu(\Pi_i).
\]
Эта величина не зависит от выбранного разбиения на клетки. Пустое множество
считается клеточным и имеет меру \(0\).

\textbf{Нижняя и верхняя меры Жордана.}
Для множества \(E\subset\R^n\) нижняя мера Жордана определяется как
\[
  \lowmu(E)=
  \sup\{\mu(A): A\text{ -- клеточное множество},\ A\subset E\},
\]
а верхняя мера Жордана как
\[
  \upmu(E)=
  \inf\{\mu(B): B\text{ -- клеточное множество},\ E\subset B\}.
\]
Если клеточного множества \(B\), содержащего \(E\), не существует, то
\(\upmu(E)=+\infty\). Для ограниченного \(E\) всегда
\(\lowmu(E)\le\upmu(E)\).

\textbf{Измеримость по Жордану.}
Множество \(E\subset\R^n\) называется измеримым по Жордану, если
\[
  \lowmu(E)=\upmu(E)\in\R.
\]
Общее значение называется мерой Жордана множества:
\[
  \mu(E)=\lowmu(E)=\upmu(E).
\]
Следовательно, всякое измеримое по Жордану множество ограничено.

\textbf{Множество меры нуль.}
Множество \(Z\subset\R^n\) имеет меру нуль, если для любого \(\eps>0\)
существует клеточное множество \(C_\eps\) такое, что
\[
  Z\subset C_\eps,\qquad \mu(C_\eps)<\eps.
\]
В частности, тогда \(\lowmu(Z)=\upmu(Z)=0\).

\section*{Теоремы}

\textbf{1. Критерий приближения клеточными множествами.}
Множество \(E\subset\R^n\) измеримо по Жордану тогда и только тогда, когда для
любого \(\eps>0\) существуют клеточные множества \(A_\eps\) и \(B_\eps\) такие,
что
\[
  A_\eps\subset E\subset B_\eps,
  \qquad
  \mu(B_\eps)-\mu(A_\eps)<\eps.
\]

\textbf{2. Критерий измеримости по границе.}
Множество \(E\subset\R^n\) измеримо по Жордану тогда и только тогда, когда
\(E\) ограничено и его граница \(\Fr E\) имеет меру нуль.

\textbf{3. Измеримость операций над множествами.}
Если множества \(E,F\subset\R^n\) измеримы по Жордану, то множества
\[
  E\cup F,\qquad E\cap F,\qquad E\setminus F
\]
также измеримы по Жордану.

\textbf{4. Конечная аддитивность меры Жордана.}
Если \(E_1,\dots,E_m\subset\R^n\) измеримы по Жордану и попарно не имеют общих
внутренних точек, то
\[
  \mu\left(\bigcup_{i=1}^m E_i\right)
  =
  \sum_{i=1}^m\mu(E_i).
\]
В частности, это верно для попарно непересекающихся измеримых множеств.

\section*{Доказательства}

\textbf{Критерий приближения клеточными множествами.}
Пусть \(E\) измеримо. Тогда \(\lowmu(E)=\upmu(E)=\mu(E)\). По определению
супремума найдется клеточное \(A_\eps\subset E\), для которого
\[
  \mu(E)-\mu(A_\eps)<\frac{\eps}{2},
\]
а по определению инфимума найдется клеточное \(B_\eps\supset E\), для которого
\[
  \mu(B_\eps)-\mu(E)<\frac{\eps}{2}.
\]
Следовательно,
\[
  \mu(B_\eps)-\mu(A_\eps)<\eps.
\]

Обратно, пусть такие \(A_\eps,B_\eps\) существуют при любом \(\eps>0\). Тогда
\[
  0\le \upmu(E)-\lowmu(E)
  \le \mu(B_\eps)-\mu(A_\eps)<\eps.
\]
Так как \(\eps\) произвольно, \(\upmu(E)=\lowmu(E)\). Кроме того, существование
клеточного \(B_\eps\supset E\) дает конечность верхней меры, поэтому \(E\)
измеримо.

\textbf{Критерий измеримости по границе.}
Пусть \(E\) измеримо. Тогда оно ограничено. Возьмем \(\eps>0\) и по
предыдущему критерию выберем клеточные множества
\[
  A_\eps\subset E\subset B_\eps,\qquad
  \mu(B_\eps)-\mu(A_\eps)<\eps.
\]
Положим
\[
  C_\eps=B_\eps\setminus \Int A_\eps.
\]
Для клеточных множеств \(B_\eps\) и \(A_\eps\) множество \(C_\eps\) клеточно.
Так как \(E\subset B_\eps\), то \(\overline E\subset B_\eps\); так как
\(A_\eps\subset E\), то \(\Int A_\eps\subset\Int E\). Поэтому
\[
  \Fr E=\overline E\setminus\Int E
  \subset B_\eps\setminus\Int A_\eps=C_\eps.
\]
Клеточные множества \(A_\eps\) и \(C_\eps\) не имеют общих внутренних точек и
\(B_\eps=A_\eps\cup C_\eps\), значит
\[
  \mu(C_\eps)=\mu(B_\eps)-\mu(A_\eps)<\eps.
\]
Итак, для любого \(\eps>0\) граница \(\Fr E\) покрывается клеточным множеством
сколь угодно малой меры, то есть \(\mu(\Fr E)=0\).

Обратно, пусть \(E\) ограничено и \(\mu(\Fr E)=0\). Зафиксируем \(\eps>0\).
Выберем клеточное множество \(C_\eps\) так, что
\[
  \Fr E\subset \Int C_\eps,\qquad \mu(C_\eps)<\eps.
\]
Такое покрытие получается из нулевой меры границы: при необходимости клеточное
покрытие немного расширяют, чтобы граница лежала во внутренности. Так как
\(E\) и \(C_\eps\) ограничены, существует клетка \(\Pi_0\), содержащая оба
множества. Множество
\[
  D_\eps=\Pi_0\setminus \Int C_\eps
\]
клеточно; представим его как объединение клеток с непересекающимися
внутренностями. Ни одна из этих клеток не пересекает \(\Fr E\), потому что
\(\Fr E\subset\Int C_\eps\). Если клетка содержит точку из \(E\) и точку вне
\(E\), то отрезок между ними, лежащий в клетке, пересекает \(\Fr E\). Поэтому
каждая клетка из разбиения \(D_\eps\) либо целиком лежит в \(E\), либо не
пересекает \(E\).

Пусть \(A_\eps\) -- объединение всех клеток из разбиения \(D_\eps\), которые
лежат в \(E\). Тогда \(A_\eps\) клеточно и \(A_\eps\subset E\). Кроме того,
любая точка \(E\), не попавшая в \(A_\eps\), лежит в \(C_\eps\), значит
\[
  E\subset A_\eps\cup C_\eps.
\]
Положим \(B_\eps=A_\eps\cup C_\eps\). Тогда
\[
  A_\eps\subset E\subset B_\eps,
  \qquad
  \mu(B_\eps)-\mu(A_\eps)\le \mu(C_\eps)<\eps.
\]
По критерию приближения \(E\) измеримо.

\textbf{Измеримость объединения, пересечения и разности.}
Сначала отметим включения для границ:
\[
  \Fr(E\cup F)\subset \Fr E\cup\Fr F,
  \qquad
  \Fr(E\cap F)\subset \Fr E\cup\Fr F,
  \qquad
  \Fr(E\setminus F)\subset \Fr E\cup\Fr F.
\]
Например, если \(x\in\Fr(E\cup F)\), то \(x\in\overline{E\cup F}\) и
\(x\notin\Int(E\cup F)\). Так как
\(\overline{E\cup F}=\overline E\cup\overline F\), точка \(x\) лежит в
\(\overline E\) или в \(\overline F\). При этом она не может быть внутренней
точкой соответствующего множества, иначе была бы внутренней точкой
\(E\cup F\). Значит \(x\in\Fr E\cup\Fr F\). Остальные два включения
доказываются тем же рассуждением через определения замыкания и внутренности.

Если \(E\) и \(F\) измеримы, то они ограничены, а их границы имеют меру нуль.
Объединение двух множеств меры нуль снова имеет меру нуль: для заданного
\(\eps>0\) покрываем каждое из них клеточным множеством меры \(<\eps/2\) и
берем объединение покрытий. Следовательно,
\[
  \Fr(E\cup F),\quad \Fr(E\cap F),\quad \Fr(E\setminus F)
\]
имеют меру нуль. Сами множества \(E\cup F\), \(E\cap F\), \(E\setminus F\)
ограничены, поэтому по критерию измеримости они измеримы по Жордану.

\textbf{Аддитивность для двух множеств.}
Пусть \(E_1,E_2\) измеримы и не имеют общих внутренних точек. Тогда
\(E_1\cup E_2\) измеримо по уже доказанному. Для произвольного \(\eps>0\)
выберем клеточные множества
\[
  A_i\subset E_i\subset B_i,\qquad
  \mu(A_i)>\mu(E_i)-\frac{\eps}{2},\qquad
  \mu(B_i)<\mu(E_i)+\frac{\eps}{2},
  \quad i=1,2.
\]
Клеточные множества \(A_1\) и \(A_2\) не имеют общих внутренних точек, поэтому
\[
  \mu(A_1\cup A_2)=\mu(A_1)+\mu(A_2).
\]
Кроме того,
\[
  A_1\cup A_2\subset E_1\cup E_2\subset B_1\cup B_2.
\]
Используя монотонность меры и субаддитивность для клеточных множеств, получаем
\[
  \mu(E_1)+\mu(E_2)-\eps
  <
  \mu(A_1)+\mu(A_2)
  =
  \mu(A_1\cup A_2)
  \le
  \mu(E_1\cup E_2),
\]
а также
\[
  \mu(E_1\cup E_2)
  \le
  \mu(B_1\cup B_2)
  \le
  \mu(B_1)+\mu(B_2)
  <
  \mu(E_1)+\mu(E_2)+\eps.
\]
Так как \(\eps>0\) произвольно,
\[
  \mu(E_1\cup E_2)=\mu(E_1)+\mu(E_2).
\]
Формула для конечного числа множеств получается индукцией.

\section*{Финальная устная версия}

В \(\R^n\) клеткой называется прямоугольный параллелепипед
\[
  \Pi=\{x: a_k\le x_k\le b_k,\ k=1,\dots,n\},
\]
его мера равна \(\prod_{k=1}^n(b_k-a_k)\). Клеточное множество -- это конечное
объединение клеток с попарно непересекающимися внутренностями; его мера равна
сумме мер клеток и не зависит от разбиения.

Для \(E\subset\R^n\) нижняя мера Жордана равна
\[
  \lowmu(E)=\sup\{\mu(A): A\text{ клеточно},\ A\subset E\},
\]
а верхняя мера Жордана равна
\[
  \upmu(E)=\inf\{\mu(B): B\text{ клеточно},\ E\subset B\}.
\]
Множество \(E\) измеримо по Жордану, если
\[
  \lowmu(E)=\upmu(E)\in\R.
\]
Общее значение называется \(\mu(E)\). Эквивалентная форма: для любого
\(\eps>0\) можно найти клеточные \(A_\eps,B_\eps\), такие что
\[
  A_\eps\subset E\subset B_\eps,\qquad
  \mu(B_\eps)-\mu(A_\eps)<\eps.
\]

Главный критерий: \(E\) измеримо по Жордану тогда и только тогда, когда оно
ограничено и \(\mu(\Fr E)=0\). Если \(E\) измеримо, берем клеточные
\(A_\eps\subset E\subset B_\eps\) с малой разностью мер. Тогда
\[
  \Fr E\subset B_\eps\setminus\Int A_\eps,
\]
а мера этого клеточного множества меньше \(\eps\). Значит граница имеет меру
нуль. Обратно, если \(\Fr E\) имеет меру нуль, покрываем ее клеточным
множеством \(C_\eps\) меры \(<\eps\), строим внутри большой клетки
\(D_\eps=\Pi_0\setminus\Int C_\eps\) и берем объединение тех клеток
\(D_\eps\), которые целиком лежат в \(E\). Получаются клеточные
\[
  A_\eps\subset E\subset A_\eps\cup C_\eps,
\]
и разность их мер меньше \(\eps\), значит \(E\) измеримо.

Если \(E\) и \(F\) измеримы, то \(E\cup F\), \(E\cap F\) и \(E\setminus F\)
ограничены, а их границы лежат в \(\Fr E\cup\Fr F\):
\[
  \Fr(E\cup F),\ \Fr(E\cap F),\ \Fr(E\setminus F)
  \subset \Fr E\cup\Fr F.
\]
Правая часть имеет меру нуль, поэтому все эти множества измеримы по критерию.

Наконец, если измеримые множества \(E_1,\dots,E_m\) попарно не имеют общих
внутренних точек, то
\[
  \mu\left(\bigcup_{i=1}^mE_i\right)=\sum_{i=1}^m\mu(E_i).
\]
Для двух множеств доказательство такое: приближаем каждое \(E_i\) изнутри и
снаружи клеточными множествами \(A_i\subset E_i\subset B_i\). Внутренние
приближения \(A_1,A_2\) не имеют общих внутренних точек, поэтому их клеточная
мера аддитивна. Снаружи используем оценку
\(\mu(B_1\cup B_2)\le\mu(B_1)+\mu(B_2)\). Получаем
\[
  \mu(E_1)+\mu(E_2)-\eps
  \le
  \mu(E_1\cup E_2)
  \le
  \mu(E_1)+\mu(E_2)+\eps,
\]
и при \(\eps\to0\) выходит равенство. Для конечного числа множеств применяется
индукция.

\section*{Источники}

Иванов Г.Е. \emph{Лекции по математическому анализу. Часть 1}:
\texttt{IvGE\_1.pdf}, стр. 200--206.

Основные роли источника: стр. 200--201 -- клетка, клеточное множество, нижняя
и верхняя меры Жордана, измеримость; стр. 202 -- критерий приближения
клеточными множествами; стр. 204--205 -- критерий измеримости по границе;
стр. 205--206 -- измеримость объединения, пересечения и разности, конечная
аддитивность меры Жордана.

Дополнительно был просмотрен \texttt{IvGE\_2.pdf}: прямого материала по
определению меры Жордана для этого билета там нет, используются только более
поздние приложения измеримости.

\end{document}
